\section{Nineteen-color constraints and higher-dimensional questions}
\label{sec:higher}

Theorem~\ref{thm:seven-lower} proves a numerical lower bound;
deciding whether nineteen colors suffice is a separate problem.
We give common necessary conditions for a hypothetical
nineteen-color line coloring and a general reduction for the
full-subspace problem.

\subsection{Deficits, projected charges and hole capacity}

Let $C_z$ be the class containing $z$, and put $s=|C_z|$.
Every other member of $C_z$ is odd, and Lemma~\ref{lem:seven-capacity}
gives $s\leq8$. The other eighteen classes have capacity seven
and contain $127-s$ vertices. If $D_1,\ldots,D_r$ are the
undersized classes, with sizes $k_i$, then
\begin{equation}\label{eq:seven-deficit}
\sum_{i=1}^r(7-k_i)=s-1.
\end{equation}
Write $m_i,o_i$ for their even and odd counts, $m=\sum_i m_i$,
and let $A,B,C$ count the saturated pure even, mixed and pure
odd heptads. Counting both parities gives
\[
7A+6B+m=63,\qquad
B+7C+\sum_i o_i+s=64,\qquad A+B+C=18-r.
\]
Using $\sum_i k_i=7r-(s-1)$, these equations become
\begin{equation}\label{eq:seven-counts}
B=m+7t,\qquad A=9-m-6t,\qquad C=9-r-t
\end{equation}
for an integer $t$ with nonnegative counts. The same formulas
apply to every characteristic size; $s$ determines the defect
patterns through \eqref{eq:seven-deficit}.

Let $\pi(v)=v+(v\cdot z)z$, $h=\sum_{v\in C_z}\pi(v)$, and
$w_i=\sum_{v\in D_i}\pi(v)$. Every saturated class has projected
sum zero. For a pure even heptad this was proved in
Theorem~\ref{thm:seven-lower}; for a pure odd heptad it follows
by summing a Lagrangian's nonzero points. In a mixed class the
six even vectors form a basis of $E$, since $J_6+I_6$ is invertible.
Their sum and the odd member's projection both pair to one with
each basis member and are therefore equal. The full projected
sum is zero, so
\begin{equation}\label{eq:seven-charge}
\sum_i w_i=h.
\end{equation}

For a quadratic refinement $q$ of $E$, an independent class with
$k$ members and $o$ odd members satisfies
\[
\sum_v q(\pi(v))
 =q\Bigl(\sum_v\pi(v)\Bigr)+\binom{k}{2}+\binom{o}{2}
 \quad\text{in }\F.
\]
The characteristic projections are pairwise orthogonal, and
contribute $q(h)$. A pure even or mixed heptad contributes one;
a pure odd heptad contributes zero. Summing over the partition
and polarizing \eqref{eq:seven-charge} gives
\begin{equation}\label{eq:seven-quadratic-charge}
\sum_{i<j} w_i\cdot w_j
 =A+B+\sum_i\left(\binom{k_i}{2}+\binom{o_i}{2}\right)
 \quad\text{in }\F.
\end{equation}
These identities allow $h=0$ and do not assume that individual
defect charges lie in a hole Lagrangian.

There is a common capacity inequality whenever the $C$ pure
odd heptads leave $d=9-C$ disjoint hole Lagrangians. Every
pure even heptad meets each isotropic side at most once. A mixed
class's odd projection forbids one whole side, so its even members
occupy at most $d-1$ holes. A pure odd defect contributes no even
hole point; a pure even defect contributes at most $d$; a defect
with both parities contributes at most $d-1$. Thus
\begin{equation}\label{eq:seven-hole-capacity}
7d\leq dA+(d-1)B+\sum_i b_d(k_i,m_i),\qquad
b_d(k,m)=
\begin{cases}
0,&m=0,\\
d,&m=k,\\
d-1,&0<m<k.
\end{cases}
\end{equation}
For $C=8$ the required completion is
Lemma~\ref{lem:eight-completion}. For $C=6,7$ it uses the finite
completion results in the accompanying package~\cite{Package}:
six or seven disjoint Lagrangians leave respectively three or two
uniquely determined hole Lagrangians. Their proof consists of
a written symplectic normalization and two independent exhaustive
enumerations; these finite premises are separate from the written
proof of Theorem~\ref{thm:seven-lower}.

For $s=4$, the outside deficit is three, giving patterns $4$,
$5+6$ and $6+6+6$. Their parity-count catalogs contain $9$, $46$
and $48$ rows. The accompanying further analysis excludes the
single-four pattern, twenty-one five-plus-six rows and fourteen three-six
rows, leaving $25$ and $34$ count candidates. They are necessary
profiles, not realized partitions. Of the remaining $59$, $57$
have $C\in\{6,7\}$ and already satisfy
\eqref{eq:seven-hole-capacity}; two have $C=5$, and no $C=8$ row remains.
The profile
$(m_5,m_6;A,B,C)=(5,4;6,2,8)$ is excluded: its four exhaustive
marked forms each have independently checked covering certificates.
For the zero-five-even-sum form, a written marked symplectic group
of order $128$, including exchange of the two mixed classes, reduces
$1\,271\,424$ necessary remaining sets to $10\,837$ representatives.
A finite failed-state proof excludes every representative and hence
closes this raw profile. The other three forms retain their separate
certificates.

The profile $(m_5,m_6;A,B,C)=(3,6;6,2,8)$ is also excluded.
Let $r$ be the sum of the five-class's three even members,
$q$ the pure-six sum, $a,d$ its five-class odd projections,
and $b,c$ the mixed projections in the hole Lagrangian $M$.
Equations~\eqref{eq:seven-charge} and
\eqref{eq:seven-quadratic-charge} give
$q+r=b+c$ and $r\cdot(b+c)=0$, while
$r\cdot a=r\cdot d=1$ puts $r$ outside $M$.
The nonzero functional $\ell(x)=r\cdot x$ therefore gives three
exhaustive role patterns: either $\ell(b)=\ell(c)=1$ and
$a+d=b+c$, or both values are zero and $a+d$ is $b+c$ or
one of $b,c$. The pure-six defect has one actual hole with
$\ell(t)=1$. Their full marked stabilizers have orders $32,32,16$;
they retain both mixed roles and transport the actual hole.
Complete necessary-root catalogs reduce to $248,474,557$
representatives. A new independently audited failed-state
certificate excludes every six-heptad cover, closing the only
remaining five-plus-six row with $C=8$.

The three-six profile
$(m_{6,1},m_{6,2},m_{6,3};A,B,C)=(2,6,6;7,0,8)$ is excluded as well.
Write $u,v$ for the even members of its two-even/four-odd defect,
$r=u+v$, and $q_1,q_2$ for the two pure-six sums. The four odd
projections exhaust the affine plane $\ell=1$ in $M$, where
$\ell=u\cdot(-)|_M=v\cdot(-)|_M$. Hence $r\in M$, $\ell(r)=1$,
and the characteristic projections are $\ker\ell\setminus\{0\}$,
so $h=0$. The charge identities give $q_1+q_2=r$ and
$q_1\cdot r=q_2\cdot r=q_1\cdot q_2=0$.
If both sums lie in $M$, exactly one is an odd projection of the
mixed defect. Moving that odd member into its corresponding pure
sextet gives the previously excluded profile $(2,6;7,1,8)$.
This recoloring step retains that earlier theorem's finite cover input.
Otherwise both sums lie outside $M$, and a written symplectic
normalization gives $\{u,v\}=\{95,111\}$, $r=48$,
$q_1=65$ and $q_2=113$, naming $q_1$ by its zero pairing with $u$.
The full marked group has order eight and fixes the two pure-six
roles individually. Its $384$ distinct necessary $49$-point sets
form $75$ orbits, with sizes $2,4,8$ occurring $12,36,27$ times.
Each pure defect may contain one hole. The four ordered hole
allocations retain $2,11,8,54$ representatives, and the cover
catalog therefore includes all $288$ even heptads: $224$ anchored
in $M$ and $64$ avoiding $M$. A new independently checked finite
certificate excludes every representative's seven-heptad cover.

The profile $(4,4,6;7,0,8)$ admits a direct recoloring reduction.
Let $D_1,D_2$ be its four-even/two-odd six-defects and $P$ its
pure even sextet. Their four odd projections lie in the hole
Lagrangian $M$. The charge-member lemma in the package names
them $a,b$ and $c,d$ so that the defect charges are $b,d$,
and their even sums are $a,c$. The characteristic projections
are the three unused points of $M\setminus\{0\}$, giving
$h=a+b+c+d$. If $q$ is the sum of $P$,
\eqref{eq:seven-charge} forces $q=a+c\in M\setminus\{0\}$.
The nonsingular sextet Gram matrix gives $p\cdot q=1$ for every
$p\in P$, so $z+q$ can be moved into $P$ to form a mixed heptad.
Its projection is a hole, so $z+q$ comes from the characteristic
class or a mixed defect. In the first case the characteristic
size becomes three, with two four-even six-defects and saturated
counts $(7,1,8)$, excluded by the package's size-three two-six
theorem. In the second case $q$ differs from $a,c$, so it is
$b$ or $d$; removing that odd member gives the size-four
five-plus-six profile $(4,4;7,1,8)$, excluded by the four-even-pair
theorem. These exact-profile transfers inherit the two theorems'
finite covering premises and require no new covering search.

The remaining profile $(4,5,5;7,0,8)$ has a further recoloring
reduction. Name the four-even defect's odd projections $a,b$,
with charge $b$ and even sum $a$, and the other defects' odd
projections $c,d$ and five-even sums $r,s$. Partition charge gives
$r+s=a$. Every five-even/one-odd six-class has a unique even
completion: $t_1=r+c$ and $t_2=s+d$ respectively. Its even sum
pairs to zero with its own even members and to one with its odd
projection, so the completion pairs to one with all six members
and lies outside $M$. The quadratic identity gives the common
functional $\ell=r\cdot(-)|_M=s\cdot(-)|_M$ the values
$\ell(a)=\ell(c)=\ell(d)=1$. Thus
$t_1+t_2=a+c+d\in M\setminus\{0\}$ and $t_1\cdot t_2=1$.
If a completion came from another defect, transferring it would
give $(3,5;7,1,8)$ or $(4,4;7,1,8)$, both previously excluded in
the package. Both completions must therefore lie in pure even
heptads. If they shared one heptad, moving both would give the
excluded profile $(5,4;6,2,8)$. Their sources are consequently
distinct. Moving both produces $(4,6,6;5,2,8)$ with two pure
sextets whose sums are $t_1,t_2$ outside $M$ and which retain
distinct actual holes. This reduces the original row to a two-hole
subbranch; the transfer alone excludes neither raw row. It retains the three earlier
theorems' historical finite premises.

This transferred two-hole branch has three exhaustive marked forms.
Put $e_1=a+c$, $e_2=a+d$, $e_3=a$; these form a basis of $M$
with the common functional taking values $(0,0,1)$.
Choosing a symplectic dual pair perpendicular to $\Span(e_3,t_1)$
normalizes $(a,c,d,t_1,t_2)$ to $(48,51,60,95,96)$.
The remaining point $b$ is $3$, $15$ or $63$, after allowing the
exchange of the paired pure-six and mixed roles. The corresponding
characteristic charges are $60$, $48$ and $0$. A marked isometry
fixes $a,b$ and either fixes both paired roles or exchanges them.
Its restriction to $M$ is the identity or the exchange of $e_1,e_2$.
There are eight symmetric-shear lifts for each admissible restriction;
the exchange fixes $b$ exactly in the last two forms. The full
marked group orders are therefore $8,16,16$. The actual ordered
hole pairs form $12,7,7$ orbits, so they cannot be collapsed into
one hole allocation. This normalization concerns the transferred
subbranch. Requiring joint disjointness of its four-even defect,
two pure sextets and two mixed even sextets gives $1708,1648,304$
distinct necessary remaining sets for $b=3,15,63$. Each has
$35$ even points, exactly five holes and a unique five-component
decomposition within its marked form. Their full-group orbits
number $224,112,21$: the first form has $21$ orbits of size four
and $203$ of size eight, the second has $18$ of size eight and
$94$ of size sixteen, and the third has four of size eight and
$17$ of size sixteen. The action transports both paired roles
and actual holes; it is not free. Each completed coloring would
partition such a remaining set into five pure even heptads.
Since a heptad meets $M$ at most once, all five must be anchored
at the remaining holes. The $357$ representatives therefore
specify a complete necessary covering problem using the $224$
anchored heptads. An independently checked failed-state certificate
excludes every representative's five-heptad cover. It has $359$
reachable states, two required branches and $357$ leaves. Each
nonempty state records a pivot, and the auditor checks every
anchored heptad contained in that state and covering its pivot.
All successors have seven fewer points and are certified failures;
the empty set is never marked failed. Induction on the remaining
size proves that no representative has a cover. This excludes
the transferred subbranch and, through the exact recoloring above,
the entire source row $(4,5,5;7,0,8)$.
The raw $(4,6,6;5,2,8)$ row retains other branches. It and
$(4,5,6;6,1,8)$ are the two three-six rows with $C=8$ at this stage.
To exclude $(4,5,6;6,1,8)$, name the four-even
defect's odd projections $a,b$, with even sum $a$ and charge $b$;
let the five-even defect have odd projection $c$ and even sum $r$,
and let the existing mixed heptad have odd projection $d$.
Write $t=r+c$ for the mixed completion and $q$ for the pure
sextet's sum and unique completion. The latter assertion follows
because the sextet's Gram matrix $I+J$ is nonsingular, its members
form a basis of $E$ and their sum pairs to one with each member.
Partition charge and the quadratic identity give
\[
 b+t+q=a+b+c+d,\qquad t+q=a+c+d\in M,\qquad t\cdot q=1.
\]
Here $A+B=7$ and the three defect terms are $0,1,1$.
Both completions lie outside $M$ and restrict to a common
functional $\ell$ with
\[
 \ell(c)=1,\qquad \ell(a)=\ell(d)=\eta\in\{0,1\}.
\]
In particular $a,c,d$ form a basis of $M$. The seven even holes
can occur only in the six pure even heptads and the pure sextet,
each at most once; thus each of these classes has one actual hole.

Moving $q$ from the four-even defect, five-even defect or existing
mixed heptad into the pure sextet would give respectively
$(3,5;7,1,8)$, $(4,4;7,1,8)$ or $(4,5,5;7,0,8)$.
Those exact rows are excluded in the package, so $q$ must lie
in a pure even heptad $H_q$. Moving $t$ from the four-even
defect or pure sextet into its recipient would give respectively
$(3,6;6,2,8)$ or $(5,4;6,2,8)$, also excluded.
Thus $t$ lies in the existing mixed heptad or a pure even heptad
$H_t$. A mixed source forces $\eta=1$; the transfer then exchanges
the mixed roles $c,d$ reversibly, retaining the same completion $t$.
If $H_t=H_q$, moving both completions would leave a pure
five-defect and give $(5,4;6,2,8)$, excluded by all four
D/H/G/Z forms. Consequently $H_q$ avoids $t$, and $H_t$,
when present, differs from $H_q$.

The pure-source transfer for $t$ produces $(4,6,6;5,2,8)$
with two distinct actual pure-six holes. One completion is in its
paired mixed class; $q$ remains in the distinct pure heptad $H_q$.
This differs from the paired mixed/mixed branch excluded above.
Both functional values and the mixed-source alternative must be
retained in the subsequent covering analysis. These source restrictions
retain the five exact-profile theorems' historical finite premises.

This source geometry has eight exhaustive ordered marked forms.
The original roles have distinct even/odd sizes $(4,2),(5,1),
(6,0),(6,1)$, so the reversible mixed-role recoloring is not
an isometry exchanging the original marked classes. Normalize
the ordered basis $(a,c,d)$ and a symplectic dual basis to
$(3,12,48,65,71,95)$. Writing $t=(x,y)$ gives
$y=(\eta,1,\eta)\ne0$. A symmetric shear
$(x,y)\mapsto(x+By,y)$ can remove $x$: for nonzero $y$, the
map $B\mapsto By$ from symmetric matrices onto $M$ is surjective,
as seen by taking $y$ to a coordinate vector. Thus
$(t,q)=(71,120)$ for $\eta=0$ and $(89,102)$ for $\eta=1$.
In each case $b$ has the four distinct positions $15,51,60,63$,
with characteristic charges $48,12,3,0$. An ordered marked
isometry fixes $a,c,d$, hence fixes $M$ pointwise. It is a symmetric
shear with $By=0$. These three independent conditions in the
six-dimensional space of symmetric matrices leave exactly eight
maps. They fix $t,q,b$ and every actual hole; this is the full
ordered marked group, with no role exchange. The eight forms
cover all charge markings and retain twelve source-case branches:
one pure-source branch for each $\eta=0$ form, and pure and mixed
branches for each $\eta=1$ form.

Jointly disjoint components $D_0,D_1,P,J,H_q,H_t$ in a pure-source
coloring have even sizes $4,5,6,6,7,7$. Their complement has
$28$ points and four holes, to be covered by four pure heptads.
For a mixed source, $D_0,D_1,P,J,H_q$ have sizes $4,5,6,6,7$,
leaving $35$ points and five holes for five pure heptads.
Every residual heptad must meet a hole, since it can meet $M$
at most once. The necessary covering catalog is therefore the
$224$ anchored even heptads, restricted to the specified remainder.
Constructing all jointly disjoint tuples from the complete local
component catalogs gives $51\,616$ ordered families in these
twelve branches. Any coloring selects one such tuple, so their
remaining sets form a complete necessary covering problem.
There are $25\,364$ form-labelled distinct remainders in $3258$
full marked-group orbits: $2604$ pure-source orbits for four
heptads and $654$ mixed-source orbits for five. The counts for
$b=15,51,60,63$ are respectively $381,384,441,170$ for
$\eta=0$ pure sources; $354,468,344,62$ for $\eta=1$ pure
sources; and $210,210,194,40$ for $\eta=1$ mixed sources.

A remainder has two, four or eight ordered component
decompositions, all retained with their actual holes. Every
group image of every tuple occurs among the decompositions
of the image remainder. Thus the full marked group transports
both roots and eligible anchored-heptad covers; one representative
of each complete remaining-set orbit suffices for a later
necessary cover test. This argument assumes neither unique
decomposition nor a free action. The earlier paired mixed/mixed
certificate does not apply to these branches.

The new finite covering certificate excludes all $3258$
representatives. It has $3262$ reachable failed states, four
required successor branches and $3258$ leaves. At each nonempty
state $R$, it records a pivot $p\in R$. The independent auditor
checks every anchored heptad $H\subseteq R$ containing $p$,
and requires every successor $R\setminus H$ to be a recorded
failure with seven fewer points. The empty set is never recorded
as failed. A leaf has no possible heptad through its pivot;
induction on remaining size therefore proves failure at every
recorded state. Every representative occurs among these states,
so no necessary four- or five-heptad cover exists. This excludes
all twelve source branches and the entire $(4,5,6;6,1,8)$ row.
Exactly one further raw row is removed, leaving $25$ five-plus-six
and $35$ three-six candidates at this stage. The sole remaining
three-six row with $C=8$, $(4,6,6;5,2,8)$, is settled below.

For that remaining row, retain the four-even defect's projections
$a,b$ and charge $b$, the two mixed projections $c,d$, and the
two pure-six sums and unique completions $u,v$. Partition and
quadratic charge now give
\[
 u+v=a+c+d\in M\setminus\{0\},\qquad u\cdot v=1,
 \qquad \ell(a)+\ell(c)+\ell(d)=1.
\]
Both completions lie outside $M$, define the same nonzero
functional $\ell$ on $M$, and $a,c,d$ form a basis.
The ordered functional patterns are $100,010,001,111$; the earlier
paired mixed/mixed branch's restriction to $111$ does not exhaust
this row. Moving either completion from the four-even defect or
the other pure sextet gives respectively the excluded
$(3,6;6,2,8)$ or $(5,4;6,2,8)$. Moving it from either mixed
heptad gives the newly excluded $(4,5,6;6,1,8)$ row.
Thus both completions have pure even sources $H_1,H_2$.
If they shared a source, moving both would give the excluded
$(5,4;6,2,8)$ row, using all four D/H/G/Z forms.
Consequently $H_1,H_2$ are distinct and each avoids the other
completion.

The five pure heptads and two pure sextets each have one actual
$M$-hole. The two sextet holes pair to one with their sums, and
each source hole pairs to one with its contained completion.
These four distinct holes therefore exhaust the four points
with $\ell=1$. The remaining three holes are exactly
$\ker\ell\setminus\{0\}$. Swapping a sextet and its pure source
by moving the completion is reversible, retaining each hole;
this recoloring is not assumed to be a geometric role symmetry.
Jointly disjoint choices of the four-even block, two mixed even
sextets, two pure sextets and their two source heptads have sizes
$4,6,6,6,6,7,7$, leaving $21$ points and the three kernel holes.
Every residual heptad must meet one of those holes, so the
eligible catalog consists of the $96$ kernel-anchored heptads,
restricted by containment. The complete construction follows.

The pure-source geometry admits sixteen exhaustive \emph{ordered}
marked forms. Retain the ordered mixed projections $c,d$ and the
ordered pure completions $u,v$, with their associated sources.
Choose a symplectic dual basis to $(a,c,d)$ and send the basis
and its dual to $(3,12,48,65,71,95)$. In these coordinates
$u=(x,y)$, where $y$ is one of $100,010,001,111$.
The pointwise-$M$ shear $(x,y)\mapsto(x+By,y)$, with $B$ symmetric,
can remove $x$: for every nonzero $y$, the map $B\mapsto By$
is onto, as seen by making $y$ a coordinate vector and choosing
the corresponding symmetric-matrix column. Thus
\[
 (u,v)=(65,126),(71,120),(95,96),(89,102),
 \qquad b\in\{15,51,60,63\}.
\]
The ordered functional values and the coordinates of $b$ in the
fixed basis are invariants, so these sixteen forms are inequivalent
under ordered marked isometries. Ordering the equal-sized roles
loses no hypothetical coloring.

Every $M$-preserving symplectic map has the form
$(x,y)\mapsto(Lx+BL^{-T}y,L^{-T}y)$ with $B$ symmetric.
Fixing $a,c,d$ forces $L=I$; fixing $u$ then gives $By=0$.
The onto map above imposes three independent conditions on six
matrix entries, yielding the full eight-element ordered marked
stabilizer for every form. It fixes $M$ pointwise, hence retains
each of the $24$ ordered allocations of four positive actual holes
to $P_1,P_2,H_1,H_2$. Completion recolorings provide no additional
ordered geometric symmetries.

Each of these four pure component catalogs has $18$ choices.
For each form, independent local enumeration gives $1152$ mutually
disjoint ordered quartets, $48$ for each actual-hole allocation.
Their full-group action has $144$ orbits of size eight; stabilizers
are checked on the tuples. These quartets occupy $26$ even points
and still require compatibility with the four-even block and both
mixed blocks. Complete seven-component tuples and their remainder
orbits must retain the ordered roles and all actual holes;
quartet freeness is not an assumption about remainder actions.
Complete joint enumeration now imposes all seven-component
disjointness conditions in every ordered form. It gives $42\,496$
ordered configurations, $10\,032$ form-labelled $21$-point remainders
and $1276$ full ordered marked-group remainder orbits.
For $b=15,51,60,63$, respectively, the orbit counts are
$44,44,184,20$ for pattern $100$;
$96,84,86,28$ for $010$;
$84,96,86,28$ for $001$; and
$136,136,100,24$ for $111$.
Every configuration retains the four actual positive holes.
All $24$ ordered hole allocations are tested; eight have an empty
joint fiber for pattern $100$ with $b=15$ or $51$, while the other
fourteen forms have all allocations nonempty.

Each remainder has four or eight ordered component decompositions,
all retained in the artifact. The joint tuple action is free and
has $5312$ orbits, but remainder orbits have sizes four or eight;
their stabilizers have orders two or one. Complete actions on all
decompositions are checked. Every hypothetical coloring normalizes
to one of these necessary tuples. Thus the complete covering
problem consists of the $1276$ representative remainders, each
requiring three heptads from the $96$ kernel-anchored rows filtered
by containment.

The finite covering certificate excludes every representative by
a one-step obstruction. For each $21$-point remainder $R$, it
records a pivot $p\in R$ such that no kernel-anchored heptad
$H\subseteq R$ contains $p$. Any partition of $R$ into three
eligible heptads would place $p$ in one of them, a contradiction.
The independent auditor reconstructs the complete roots and groups,
enumerates all even heptads, and checks this absence for every pivot.
All $1276$ recorded failed states are leaves; there are no required
successors. The full marked groups preserve the kernel holes and
eligible heptads, so failure transports to every remainder in each
orbit. Since the pure-source reduction exhausts the raw row, this
excludes the entire $(4,6,6;5,2,8)$ profile. Exact count reconstruction
removes that one row, leaving $25/34$ candidates, with $57$ having
$C=6$ or $7$ and two having $C=5$. No size-four $C=8$ count
candidate remains in either list.
This is a partial case analysis.
The package records its written reductions and finite dependencies;
it supplies no uniform exclusion of all remaining profiles.

The further analysis also excludes $s<4$. Characteristic sizes
$4$ through $8$ remain open, and larger $s$ introduce outside
deficits $4$ through $7$. Closing the $59$ remaining size-four count candidates
would therefore not finish the nineteen-color problem. Excluding
every nineteen-color line coloring would improve the lower bound
to twenty, without establishing a matching upper bound. Conversely,
a nineteen-color line witness would still require an extension
check before it proved a full-subspace upper bound.

\subsection{A dimension reduction}

\begin{proposition}\label{prop:dimension-deletion}
Let $K$ be a fixed clique in $\OG_n^*$ of size at most $k$,
precolored with distinct labels from $k$ colors. If $N(n-d)<k$,
every uncolored $d$-dimensional vertex can be deleted by the
degree-versus-list rule. The residual instance is $k$-colorable
if and only if the full graph is $k$-colorable.
\end{proposition}
\begin{proof}
Every neighbor of a $d$-space $U$ is a nonzero subspace of
$U^\perp$, so $\deg U\leq N(n-d)$. If $f$ fixed clique vertices
are adjacent to $U$, its residual degree is at most $N(n-d)-f$
and its list has size $k-f$. A color is therefore available when
$U$ is restored. Apply this argument in reverse deletion order.
\end{proof}

For $n=7$, the nonzero subspaces of $T=\Span(3,12,48)$ together
with $\Span(64)$ give a $16$-clique. Since $N(3)=15$, vertices
of dimension at least four can be deleted in any coloring problem
with at least sixteen labels and these fixed clique colors.
The lower bound in Section~\ref{sec:seven} rules out sixteen,
seventeen and eighteen colors; the reduction remains available
for the nineteen-color extension problem.

The exact chromatic numbers for $n\geq7$ and the question of
infinitely many dimensions with $\chi>\omega$ remain open.
Within dimension six, a geometric construction of an optimal
full-subspace coloring is a separate problem. The negative
results in Section~\ref{sec:geometry} show that such a construction
must use enough placement information to break the stated symmetry.

\begin{problem}\label{prob:structural}
Can one specify auxiliary geometric data and a nontrivial proper
subgroup $H<\operatorname{Stab}_{O_6(\F)}(T)$, and construct a
proper $15$-coloring $c$ from those data with $c(hU)=h\cdot c(U)$
for $h\in H$, where labels are the nonzero subspaces of $T$?
The construction should prove availability and separation without
looking up the archived coloring certificate.
\end{problem}

The checked twelve-coloring already determines $\chi(\Gamma_6)$;
passing to all nonzero subspaces increases the chromatic number
by exactly three.
